% Generated draft: requires mathematical and source review.
\documentclass[11pt]{article}
\usepackage[margin=25mm]{geometry}
\usepackage{fontspec}
\setmainfont{Latin Modern Roman}
\usepackage{amsmath,amssymb,mathtools,unicode-math}
\setmathfont{Latin Modern Math}
\usepackage{enumitem,needspace,xurl,hyperref}
\setlength{\parindent}{0pt}
\setlength{\parskip}{5pt}
\setlength{\emergencystretch}{3em}
\newcommand{\N}{\mathbb{N}}
\newcommand{\Z}{\mathbb{Z}}
\newcommand{\Q}{\mathbb{Q}}
\newcommand{\R}{\mathbb{R}}
\newcommand{\C}{\mathbb{C}}
\newcommand{\F}{\mathbb{F}}
\newcommand{\A}{\mathbb{A}}
\newcommand{\PP}{\mathbb{P}}
\newcommand{\E}{\mathbb{E}}
\newcommand{\eps}{\varepsilon}
\newcommand{\dd}{\,\mathrm{d}}
\begin{document}
\section*{UnsolvedMath ID: 30003104}
\textbf{Source problem code:} catalogue entry
\subsection{Definitions and mathematical statement}
Let \(G\) be a quasi-split group and let \(\varphi\) be a discrete Langlands parameter in the setting intended by the source. The base field, the precise class of groups, the parameter domain, and the meaning of discreteness are not specified in the supplied record. Consequently, the formulation below records the conjectural target schematically; it is not a fully specified conjecture.

In a local Langlands formulation, write \(\widehat G\) for the complex dual group and \({}^{L}G\) for the corresponding \(L\)-group. The full parameter centralizer would be
\[
S_{\varphi}=\operatorname{Cent}_{\widehat G}\bigl(\operatorname{im}\varphi\bigr),
\]
where centralization is computed in \({}^{L}G\). Let \(\mathcal I_{\mathrm{ext}}(G)\) denote the intended indexing set of extended pure inner twists of \(G\), and write \(G_{\xi}\) for the group underlying a twist \(\xi\). Neither the definition of these twists nor their equivalence relation is supplied. Let \(\Pi_{\varphi}(G_{\xi})\) denote the conjectural \(L\)-packet associated with \(\varphi\) on that twist.

The supplied assertion asks whether, for every such discrete parameter \(\varphi\), there is a parametrization of the schematic form
\[
\coprod_{\xi\in\mathcal I_{\mathrm{ext}}(G)}
\Pi_{\varphi}(G_{\xi})
\;\longleftrightarrow\;
\operatorname{Irr}(S_{\varphi}),
\]
compatible with the source's stated Kottwitz bijection between basic isocrystal classes and the character group of the dual center. Here the disjoint union retains the twist label. The meaning of \(\operatorname{Irr}(S_{\varphi})\) is unresolved: the record does not specify a representation category, whether a component group or another modification of the centralizer is intended, or the central-character conditions that assign representations to individual twists. It also does not specify the field-dependent definition of basic isocrystal classes, the relevant Galois invariants of the dual center, or the exact compatibility map. Thus these features cannot be made into a precise bijection or commutative diagram from the supplied evidence alone.

The requested parametrization must additionally have the following properties. On the quasi-split form, \(\Pi_{\varphi}(G)\) is to contain exactly one Whittaker-generic constituent, and the packets and their parametrizations are to satisfy the expected endoscopic identities. Whittaker genericity conventionally requires a Whittaker datum: in a setting with a maximal unipotent subgroup \(U\) and a nondegenerate character \(\psi\), it means that the representation \(\pi\) admits a nonzero Whittaker functional, expressed as \(\operatorname{Hom}_{U}(\pi,\psi)\ne0\). No Whittaker datum is supplied. Moreover, discreteness is not explicitly accompanied by boundedness or temperedness, so the scope of the genericity assertion requires clarification. No endoscopic data, transfer factors, character-distribution conventions, or explicit endoscopic equalities are given. The endoscopic clause therefore remains an unresolved requirement rather than a specified identity.
\subsection{Short English statement}
For each discrete Langlands parameter, the source proposes a joint parametrization of packets across all extended pure inner twists, compatible with Kottwitz's classification, with a unique Whittaker-generic member on the quasi-split form and suitable endoscopic identities. The field, centralizer representation category, twist indexing, genericity hypotheses, and identities require clarification before this becomes a precise mathematical conjecture.
\subsection{Sources}
\begin{itemize}
\item Local UnsolvedMath registry: permanent ID \(30003104\), exact source problem code \texttt{catalogue entry}, title ``Langlands Packets Across Extended Pure Inner Forms,'' in the collection ``Oberwolfach Reports Open Problems.'' No individual proposer is supplied; the proposed year is recorded as 2016.
\item Supplied report attribution: \emph{Arbeitsgemeinschaft: The Geometric Langlands Conjecture (2017)}, Oberwolfach Reports, pp. 1027--1098, DOI \texttt{10.4171/owr/2016/20}; \url{https://doi.org/10.4171/owr/2016/20}. The supplied literature-source label is ``OWR 2016 report.'' The differing year indications are retained without resolution. The background describes Oberwolfach Reports as documenting workshops and their open-problem discussions.
\item Supplied literature source, labeled ``Kaletha'': \url{https://annals.math.princeton.edu/2016/184-2/p06}. No title or specific supporting passage is supplied.
\item Supplied literature source, labeled ``Kaletha'': \url{https://arxiv.org/abs/2008.04472}. No title or specific supporting passage is supplied.
\item Supplied literature assessment dated exactly \texttt{2026-08-21}: the registry labels the problem ``partially solved'' and states that the extended-pure-inner-form refinement has major case-by-case results and precise conjectural formulations, while the general packet parametrization and endoscopic identities are not fully proved. This is supplied dated curation derived from the OWR material, not an independently verified literature assessment. The background expressly calls for independent verification of time-sensitive status and cited claims.
\item The registry's statement-status label is \texttt{corrected\_verified}; its accompanying explanation concerns repairs to notation, typography, and sentence structure without changing mathematical content. This label does not establish mathematical correctness or verification of the linked sources.
\end{itemize}
\end{document}
